\documentclass[12pt,a4paper,oneside]{book}
\usepackage[utf8]{inputenc}
\usepackage[T1]{fontenc}
\usepackage[italian]{babel}
\usepackage{graphicx} 
\usepackage{amsmath}
\usepackage{amsfonts}
\usepackage{amssymb}
\usepackage{textalpha}
\usepackage{booktabs}
\usepackage{geometry} 
\usepackage[backend=biber, style=numeric]{biblatex}
\usepackage{csquotes}
\usepackage{tcolorbox}



\newcommand{\aegistrust}{$\Lambda\Xi$GIST\reflectbox{R}UST}
\newcommand{\trustlevel}[1]{\delta_{#1} \rightarrow \infty}
\newcommand{\pred}[1]{P_{#1}(\lambda)}
\newcommand{\deepcheck}[1]{\Gamma_{df}(#1)}

\newenvironment{industrialNote}
{\begin{tcolorbox}[colback=blue!5,colframe=blue!75!black,title=Specifiche Industriali]}
{\end{tcolorbox}}

\begin{document}

\begin{titlepage}
    \centering
    \vspace{1cm}
    {\Large \textsc{Università degli Studi di Verona}} \\
    \vspace{0.5cm}
    {\large Dipartimento di Informatica} \\ 
    \vspace{0.5cm}
    {\large Elaborato Latex} \\
    
    \vspace{4cm}
    
    {\LARGE \textbf{Sviluppo e Implementazione del Software}} \\
    \vspace{0.5cm}
    {\Huge \textbf{\aegistrust}} \\
    \vspace{1cm}
    {\large \textit{Protezione Deepfake, Intelligence OSINT,Analisi Biometrica Multimediale e Aziendale}}
    
    \vspace{4cm}
    
    \begin{minipage}[t]{0.47\textwidth}
        \raggedright
        \textbf{Relatore:}\\
        Chiar.mo Prof.\\
        Enrico Gregorio
    \end{minipage}
    \hfill
    \begin{minipage}[t]{0.47\textwidth}
        \raggedleft
        \textbf{Candidato:}\\
        L$\Xi$GION\\
    \end{minipage}
    
    \vfill
    
    \rule{0.8\textwidth}{0.4pt}\\
    \vspace{0.5cm}
    {\large Anno Accademico 2025-2026}
\end{titlepage}

\tableofcontents

\chapter{Introduzione al Progetto \aegistrust}

\section{Visione d'Insieme}
Il software \aegistrust{} nasce come risposta alla crescente necessità di sicurezza proattiva. A differenza dei sistemi reattivi tradizionali, \aegistrust{} implementa un framework di fiducia dinamica basato su tre pilastri fondamentali:
\begin{itemize}
    \item Scansione continua di fonti aperte (OSINT) per la mitigazione del rischio.
    \item Analisi euristica dei comportamenti di rete mediante intelligenza artificiale.
    \item Protezione avanzata contro la manipolazione biometrica, media sintetici e furto d'identità vocale.
\end{itemize}

\section{Architettura Modulare Enterprise}
L'efficienza di \aegistrust{} è garantita dalla sua struttura a microservizi indipendenti, che permette una scalabilità orizzontale del sistema:
\begin{enumerate}
    \item \textbf{Modulo Alpha}: Gestione delle identità e dei permessi di accesso.
    \item \textbf{Modulo Sigma}: Motore di ricerca OSINT e data mining su larga scala.
    \item \textbf{Modulo Tau}: Analisi predittiva e reporting per il management.
    \item \textbf{Modulo Omega}: Layer di protezione meeting online (Zoom, Teams) con funzionalità biometriche.
\end{enumerate}

\chapter{Ricerca OSINT e Analisi Database}

\section{Integrazione Database Aziendali e Pubblici}
\aegistrust{} opera un cross-check istantaneo tra database pubblici e archivi aziendali privati. Il sistema interroga i registri delle Camere di Commercio per validare l'esistenza legale delle entità durante le trattative commerciali.

\begin{industrialNote}
Il software è progettato per la protezione della proprietà intellettuale. Ogni interrogazione ai database segue protocolli di crittografia asimmetrica per evitare il leak di informazioni sensibili durante le fasi di Due Diligence.
\end{industrialNote}

\section{Metodologie OSINT Avanzate}
Il modulo OSINT di \aegistrust{} esegue una ricerca profonda su repository pubblici e leak database per identificare vulnerabilità prima che vengano sfruttate da attori malevoli. La metrica di esposizione $E$ viene calcolata come:
\[ E = \frac{\sum_{i=1}^{n} v_i}{T} \]
dove $v_i$ rappresenta il peso della singola criticità rilevata.

\chapter{Protezione Deepfake e Analisi Biometrica}

\section{Validazione Biometrica rPPG}
Uno dei componenti più innovativi di \aegistrust{} è il motore di rilevamento Deepfake applicato ai meeting online. Il software analizza il segnale video in tempo reale alla ricerca di artefatti di generazione neurale (GAN).

Il sistema utilizza la fotopletismografia remota (rPPG) per rilevare il battito cardiaco facciale, parametro assente nei media sintetici, garantendo l'autenticità dei partecipanti durante videoconferenze critiche su piattaforme come Zoom e Microsoft Teams \cite{deepfake2026}.

\section{Analisi Comportamentale e Liveness Detection}
Oltre all'analisi rPPG, l'integrità del soggetto viene monitorata attraverso parametri fisiologici non riproducibili fedelmente dalle attuali IA generative:
\begin{itemize}
    \item \textbf{Spontaneous Blinking Frequency (SBF)}: Monitoraggio del battito delle palpebre per rilevare pattern innaturali (troppo frequenti o assenti).
    \item \textbf{Audio-Visual Coherence}: Verifica della correlazione tra l'ampiezza del volume audio e l'apertura fisica della bocca (pixel aperture).
\end{itemize}

Il calcolo della sincronia audio-video viene effettuato tramite il seguente integrale:
\begin{equation}
    \Psi_{sync} = \int_{t_0}^{t_1} | \text{Volume}(t) - \text{AperturaBocca}(t) | \, dt
\end{equation}

Se lo scostamento $\Psi_{sync}$ supera la soglia di tolleranza, il sistema segnala un tentativo di injection video o doppiaggio sintetico.

\section{Identificazione Vocale e Analisi Fonetica}
\aegistrust{} integra un motore di scansione vocale autonoma. Il software analizza le singole parole pronunciate durante il meeting, estraendo l'impronta fonetica unica dell'interlocutore per confermarne l'identità rispetto ai campioni pre-caricati nel database aziendale.

\begin{industrialNote}
L'analisi vocale avviene rigorosamente in locale (On-Premise) per garantire che i dati biometrici della voce non lascino mai il perimetro di sicurezza dell'azienda, prevenendo furti di identità su larga scala.
\end{industrialNote}

\section{Algoritmo di Sintesi Deepfake Check}
La decisione finale sulla veridicità del flusso multimediale unifica tutti i parametri biometrici (rPPG, Artefatti FFT e Sincronia) nella seguente funzione logistica aggiornata:

\begin{equation}
    \deepcheck{S} = \frac{1}{1 + e^{-(\alpha \cdot \text{rPPG} + \beta \cdot \text{FFT} + \gamma \cdot \Psi_{sync})}}
\end{equation}

\section{Modelli Predittivi Aziendali}
La predittività dei prodotti aziendali rimane una funzione core. Utilizzando modelli statistici, il software calcola la probabilità di successo $P$:

\begin{equation}
    \Omega= \oint_{\partial \Sigma} \vec{F} \cdot d\vec{s} + \sqrt{\frac{\int_{0}^{\infty} \pred{t} e^{-kt} dt}{n!}}
\end{equation}

\chapter{Specifiche Tecniche e Risultati}

\section{Analisi delle Performance Industriali}
I test eseguiti su infrastrutture aziendali mostrano un abbattimento dei falsi positivi nelle ricerche OSINT.

\begin{table}[ht]
\centering
\begin{tabular}{lccc}
\toprule
Modulo & Target Database & Tempo Risposta (ms) & Affidabilità \\
\midrule
OSINT Public & Camere Commercio & 120 & 98.2\% \\
Deepfake Shield & Zoom/Teams Feed & 45 & 99.1\% \\
Leak Scanner & Dark Web Repos & 210 & 95.5\% \\
\bottomrule
\end{tabular}
\caption{Performance dei moduli core di \aegistrust{} in ambito Enterprise.}
\end{table}

\section{Validazione dei Parametri Biometrici}
I test specifici sui moduli di riconoscimento vocale e video confermano l'efficacia contro attacchi di spoofing.

\begin{table}[ht]
\centering
\begin{tabular}{lccc}
\toprule
Parametro & Metodo Analisi & Soglia Errore & Risultato \aegistrust \\
\midrule
Identità Vocale & Analisi Fonetica & < 0.5\% & \textbf{Conforme} \\
Deepfake Video & Blinking \& rPPG & < 0.2\% & \textbf{Conforme} \\
Sync Audio-Video & Apertura/Decibel & < 1.0\% & \textbf{Conforme} \\
\bottomrule
\end{tabular}
\caption{Validazione dei parametri biometrici multimodali.}
\end{table}

\section{Integrazione Grafica}
Il flusso operativo tra il software e i database esterni è illustrato nel diagramma di sistema.

\begin{figure}[ht]
\centering
\includegraphics[width=0.6\textwidth]{grafico.pnj.png}
\caption{Diagramma di interazione tra \aegistrust, Zoom e i database pubblici.}
\end{figure}

\printbibliography

\end{document}